\section{Section I: Short-Answer Questions}
\label{sec:one}

\noindent\textbf{Questions answered:} \SectionIChoice. Each answer is at most 300 words. Equations
and tables are excluded from the count.

% ------------------------------------------------------------------
\subsection{Question 1: Linear classifiers and the perceptron}
\wordlimit{300}
Setting: a cooperative bank flags a loan as risky when $s = \vect{w}\cdot\vect{x} + b \geq 0$.

\question{1.1}{2 marks}{Why might a random linear classifier look acceptable on a handful of loans yet be unreliable? Name one thing you would do to judge it fairly.}
A random $\vect{w}$ and $b$ have learned nothing from the data, but on a handful of loans they
can look good by luck. With few points, a random boundary may happen to put most of them on the
correct side. If most sampled loans are safe, a boundary that calls almost everything safe also
scores a high accuracy. Such results come from chance, not from a real pattern, so they will
not repeat on new loans.

To judge it fairly, I would test it on a large held-out set of loans that were not used to
choose $\vect{w}$ and $b$, compare it with a majority-class baseline, and report recall and
precision, not only accuracy.

\question{1.2}{2 marks}{What does the perceptron do to $\vect{w}$ and $b$ on a misclassified loan, and why does this move the boundary helpfully?}
The perceptron changes $\vect{w}$ and $b$ only when it gets a loan wrong. With labels
$y,\hat y \in \{0,1\}$ and learning rate $\eta>0$:
\[
  \vect{w} \leftarrow \vect{w} + \eta\,(y-\hat y)\,\vect{x}, \qquad b \leftarrow b + \eta\,(y-\hat y).
\]
If a risky loan ($y=1$) is wrongly called safe ($\hat y=0$), it adds $\eta\vect{x}$ to
$\vect{w}$ and $\eta$ to $b$. If a safe loan is wrongly called risky, it subtracts them.

This helps because the new score of that loan moves towards the correct side:
\[
  s_{\text{new}} = (\vect{w} + \eta\vect{x})\cdot\vect{x} + (b+\eta)
               = s + \eta\left(\lVert\vect{x}\rVert^2 + 1\right).
\]
Since $\lVert\vect{x}\rVert^2 + 1 > 0$, a missed risky loan's score always rises and a wrongly
flagged safe loan's score always falls. The boundary $\vect{w}\cdot\vect{x}+b=0$ shifts
towards putting this loan on its correct side.

\question{1.3}{2 marks}{Describe a situation in which a perceptron would never stop updating, and explain why.}
This happens when the training data are \emph{not linearly separable}: no straight boundary can
put every risky loan on one side and every safe loan on the other. For example, two loans with
identical features where one defaulted and one was repaid.

The perceptron stops only after a full pass with zero mistakes. If the data cannot be separated,
some loan is always misclassified, so every pass triggers another update and the weights keep
moving for ever. In practice we cap the number of passes, or use a method that accepts errors,
such as logistic regression or a soft-margin SVM.

% ------------------------------------------------------------------
\subsection{Question 2: Probabilities, classes and thresholds}
\wordlimit{300}
Setting: logistic regression predicts the probability that a patient misses a follow-up visit.

\question{2.1}{2 marks}{Distinguish a predicted probability from a predicted class, and explain the role of the classification threshold.}
The \emph{predicted probability} $\hat p$ is a number between 0 and 1 giving the model's
estimate of how likely the patient is to miss the visit, for example $\hat p = 0.73$. It comes
from the sigmoid function, $\hat p = 1/(1+e^{-z})$.

The \emph{predicted class} $\hat y$ is a yes or no decision: $1$ means ``will miss'' and $0$
means ``will attend''.

The threshold $t$ turns the probability into the class:
\[
  \hat y =
  \begin{cases}
    1 & \text{if } \hat p \geq t,\\
    0 & \text{if } \hat p < t.
  \end{cases}
\]
The threshold is a policy choice, not learned from the data. A lower $t$ flags more patients,
missing fewer real cases but raising more false alarms. A higher $t$ flags fewer.

\question{2.2}{2 marks}{Patients P ($\hat p = 0.51$) and Q ($\hat p = 0.97$) are both flagged at threshold $0.5$. Is it reasonable to treat them identically? Defend a position.}
\textbf{Position: no, they should not be treated identically.}

Both cross the threshold, but their risks differ greatly. About half of patients like P would
miss the visit, while almost all patients like Q would. The binary flag hides this.

Staff time is limited. Giving both the same costly action, such as a home visit, spreads effort
thinly. A tiered response is better: Q gets a phone call and a home visit, while P gets a
low-cost SMS reminder.

If the only action is cheap, such as one SMS for everyone flagged, identical treatment would
also be reasonable.

\question{2.3}{2 marks}{Describe one setting for a threshold below $0.5$ and one for a threshold above $0.5$, with reasons.}
\textbf{Below 0.5 (for example $t = 0.3$).} Use this when missing a real case is costly and the
action is cheap. For tuberculosis patients, a missed follow-up can cause treatment failure, while
an SMS reminder costs almost nothing. A low threshold catches more patients who will miss
(higher recall), and the extra false alarms only cost a few messages.

\textbf{Above 0.5 (for example $t = 0.8$).} Use this when the action is expensive or limited.
With only a few health workers for home visits, a high threshold sends them only to patients
very likely to miss (higher precision), so their time is not wasted.

% ------------------------------------------------------------------
\subsection{Question 3: Evaluation under class imbalance}
\wordlimit{300}
Setting: about 2 in 100 households are eligible but unregistered, and a vendor claims ``98\% accuracy''.

\question{3.1}{2 marks}{Why does 98\% accuracy not show the model is useful? Use a trivial classifier in your explanation.}
Take a trivial classifier that predicts ``not eligible'' for every household. In a group of
100 households it is right for the 98 that are not eligible and wrong for the 2 that are:
\[
  \text{Accuracy} = \frac{\text{TP}+\text{TN}}{\text{Total}} = \frac{0 + 98}{100} = 0.98.
\]
So 98\% accuracy is what you get by doing nothing. This classifier finds none of the
eligible households:
\[
  \text{Recall} = \frac{\text{TP}}{\text{TP}+\text{FN}} = \frac{0}{0+2} = 0.
\]
Because eligible households are so rare, accuracy is dominated by the easy majority and says
nothing about whether the model finds the households the scheme is looking for.

\question{3.2}{2 marks}{Which two metrics would you request instead, and which question does each answer in this scheme?}
\textbf{Recall}, $\text{TP}/(\text{TP}+\text{FN})$. It answers: ``Of all the eligible
households that are not yet registered, how many does the model find?'' Low recall means
families who need food support are left out.

\textbf{Precision}, $\text{TP}/(\text{TP}+\text{FP})$. It answers: ``Of the households the
model flags, how many are really eligible?'' Low precision means field staff waste visits
verifying households that are not eligible.

\question{3.3}{2 marks}{Explain in plain language what a ROC-AUC of $0.5$ means.}
ROC-AUC measures how well the model ranks households. If we pick one eligible household and
one non-eligible household at random, the AUC is the chance that the model gives the eligible
one the higher score.

An AUC of 0.5 means this happens only half the time, which is the same as tossing a coin. The
model cannot tell eligible households apart from the rest. Changing the threshold will not
help, because the scores themselves carry no useful information.

% ------------------------------------------------------------------
\subsection{Question 4: KNN and decision trees}
\wordlimit{300}
Setting: predicting seed adoption from land size, distance to market, rainfall and schooling.

\question{4.1}{2 marks}{If rainfall were recorded in metres instead of millimetres, what happens to KNN and decision-tree predictions, and why do they differ?}
\textbf{KNN: predictions change.} KNN finds neighbours using a distance over all features:
\[
  d(\vect{u},\vect{v}) = \sqrt{\textstyle\sum_j (u_j - v_j)^2}.
\]
In millimetres, a rainfall gap such as $1200 - 800 = 400$ dominates the gaps in land size or
schooling. In metres it becomes $1.2 - 0.8 = 0.4$ and almost disappears. Different farmers
become the nearest neighbours, so the predictions change.

\textbf{Decision tree: predictions stay the same.} A tree splits on one feature at a time, such
as ``rainfall $\leq 900$ mm'', which simply becomes ``rainfall $\leq 0.9$ m''. The same farmers
fall on each side, so the tree and its predictions are unchanged.

\textbf{Why they differ:} KNN mixes all features into one distance, so units matter. A tree
compares each feature only with itself, so only the order of values matters.

\question{4.2}{2 marks}{How does overfitting show up in KNN (as $K$ changes) and in a decision tree (as depth changes)?}
\textbf{KNN:} overfitting appears when $K$ is very small. With $K = 1$ each prediction copies
the single nearest farmer, noise included, so training accuracy is perfect but accuracy on new
farmers is lower. A larger $K$ smooths the predictions, and a very large $K$ underfits by
predicting the majority class.

\textbf{Decision tree:} overfitting appears when the tree is very deep. Leaves end up holding
one or a few farmers, so the tree memorises the training data. Training accuracy keeps rising
while validation accuracy falls. A very shallow tree underfits.

In both, overfitting shows as a large gap between training and validation performance.

\question{4.3}{2 marks}{For a panchayat committee, which model would you choose and what would you give up?}
\textbf{Choice: a decision tree.}

A tree can be shown as a few if-then rules, for example: ``If land size is above 2 acres and the
market is within 5 km, the farmer is likely to adopt.'' Committee members can follow, question
and check each rule against local knowledge. KNN can only say ``similar farmers adopted'',
which hides which factors matter.

\textbf{What I give up:} a readable tree must be shallow, so it may be less accurate than KNN.
Trees also draw step-like boundaries and can change a lot with small data changes.

% ------------------------------------------------------------------
\subsection{Question 5: Margins and support vectors}
\wordlimit{300}
Setting: one feature $x$ (years of schooling). Class $-1$ at $x\in\{2,5\}$ and class $+1$ at $x\in\{8,10\}$. Hard margin.

\question{5.1}{2 marks}{State the maximum-margin boundary, the total margin width, and the support vectors.}
The closest points from the two classes are $x=5$ (class $-1$) and $x=8$ (class $+1$). The
maximum-margin boundary lies halfway between them:
\[
  x^{*} = \frac{5+8}{2} = 6.5, \qquad \text{total margin width} = 8 - 5 = 3.
\]
The support vectors are $x = 5$ and $x = 8$. The rule is: predict $+1$ if $x > 6.5$,
otherwise $-1$.

\question{5.2}{2 marks}{A class $-1$ point is added at $x=7$. Give the new boundary and support vectors. What does this show?}
The closest opposite pair is now $x = 7$ (class $-1$) and $x = 8$ (class $+1$):
\[
  x^{*} = \frac{7+8}{2} = 7.5, \qquad \text{total margin width} = 8 - 7 = 1.
\]
The new support vectors are $x = 7$ and $x = 8$. The point $x = 5$ is no longer a support
vector.

This shows that the boundary depends only on the points closest to it (the support vectors).
Points far from the boundary, such as $x = 2$ and $x = 10$, can move or be removed without
changing the boundary. A single new point near the boundary can change it a lot.

\question{5.3}{2 marks}{Instead, a class $-1$ point is added at $x=9$ (non-separable). How do small and large $C$ treat it, and which do you prefer?}
The point $x = 9$ lies between the class $+1$ points at 8 and 10, so no single cut-off can
classify every point correctly. A soft-margin SVM allows violations and charges a penalty of
$C$ for each unit of violation.

\textbf{Small $C$:} violations are cheap. The SVM keeps a wide margin decided by the main
groups and simply accepts $x = 9$ as an error, treating it like an outlier.

\textbf{Large $C$:} violations are expensive, so the SVM tries hard to reduce the error at
$x = 9$ and becomes sensitive to this one household. Here, no cut-off can put 9 on the
$-1$ side without also misclassifying the $+1$ point at 8, so a large $C$ cannot remove the
error. It only makes the solution depend more on a single unusual point, which risks
overfitting.

\textbf{Preference: a small $C$.} One unusual household should not decide the rule for
everyone, and a wider margin usually generalises better to new households. A larger $C$
would make sense only if we knew that households like $x = 9$ form a real, important group.
